
\documentclass[11pt,openany]{ctexbook}

\usepackage[a4paper,margin=2.45cm]{geometry}
\usepackage{amsmath,amssymb,amsfonts,bm,mathtools}
\usepackage{booktabs}
\usepackage{enumitem}
\usepackage{hyperref}
\usepackage{xcolor}
\usepackage{microtype}

\hypersetup{
  colorlinks=true,
  linkcolor=blue,
  urlcolor=blue,
  citecolor=blue,
  pdftitle={从粒子到计算网格},
  pdfauthor={}
}

\setcounter{tocdepth}{0} % 目录显示 part 和 chapter
\setlength{\parindent}{2em}
\setlength{\parskip}{0.35em}

\newcommand{\dd}{\,\mathrm d}
\newcommand{\ee}{\mathrm e}
\newcommand{\ii}{\mathrm i}
\newcommand{\R}{\mathbb R}
\newcommand{\Order}{\mathcal O}
\newcommand{\Lop}{\mathcal L}
\newcommand{\Nop}{\mathcal N}
\newcommand{\Ksp}{\mathcal K}
\newcommand{\avg}[1]{\left\langle #1\right\rangle}

\title{从粒子到计算网格\\
\large 面向数学背景读者的计算力学入门}
\author{}
\date{}

\begin{document}

\maketitle
\tableofcontents

\chapter*{前言}
\addcontentsline{toc}{chapter}{前言}
计算力学不是“把物理公式放进计算机”，而是一条连续的信息压缩链：
\[
\text{多体动力学}
\longrightarrow
\text{统计描述}
\longrightarrow
\text{连续介质模型}
\longrightarrow
\text{离散方程}
\longrightarrow
\text{矩阵系统}
\longrightarrow
\text{工程计算}.
\]

真实系统、数学模型、离散模型和代数系统是四个不同对象。物理建模决定忽略什么，数值离散决定近似什么，线性代数决定如何把离散问题真正算出来，网格则负责把几何、边界条件和局部尺度连接到离散系统中。

本文面向学过数学分析、实分析、复分析、常微分方程、偏微分方程和数值计算，但缺少传统物理课程背景的读者。叙述以物理建模为主，公式用于支撑概念，而不追求公理化完备。

% =========================================================
\part{理论框架}

\chapter{从多粒子动力学到动理学方程}

设系统中有 \(N\) 个粒子，第 \(i\) 个粒子的位置和动量为
\[
\bm x_i,\qquad \bm p_i.
\]
整个系统的状态是 \(6N\) 维相空间中的一点：
\[
\Gamma=
(\bm x_1,\dots,\bm x_N,\bm p_1,\dots,\bm p_N).
\]

若 Hamiltonian 为 \(H\)，粒子运动满足
\[
\dot{\bm x}_i=\frac{\partial H}{\partial \bm p_i},
\qquad
\dot{\bm p}_i=-\frac{\partial H}{\partial \bm x_i}.
\]

对宏观物质逐粒子积分并不可行，因此改用相空间概率密度
\[
\rho_N(\Gamma,t).
\]
概率守恒给出高维连续性方程
\[
\frac{\partial \rho_N}{\partial t}
+
\nabla_\Gamma\cdot(\rho_N\dot\Gamma)=0.
\]
Hamilton 系统满足
\[
\nabla_\Gamma\cdot\dot\Gamma=0,
\]
故有
\[
\frac{\partial \rho_N}{\partial t}
+
\dot\Gamma\cdot\nabla_\Gamma\rho_N=0,
\]
或
\[
\frac{\partial \rho_N}{\partial t}+\{\rho_N,H\}=0.
\]

刘维尔方程表达的是相空间中的概率云团随微观动力学被搬运，但相空间流不压缩概率体积。它仍然是 \(6N\) 维方程，真正的降维来自边缘化。

定义单粒子分布函数
\[
f(\bm x,\bm v,t),
\]
使
\[
f(\bm x,\bm v,t)\dd\bm x\dd\bm v
\]
表示位置和速度落在相应微元中的粒子期望数。

从 \(N\) 粒子方程对其余变量积分，会出现
\[
f_1\longrightarrow f_2\longrightarrow f_3\longrightarrow\cdots,
\]
即 BBGKY 层级。单粒子方程依赖双粒子关联，双粒子方程又依赖三粒子关联。要得到封闭模型，必须对关联作假设。

一般动理学方程写成
\[
\frac{\partial f}{\partial t}
+\bm v\cdot\nabla_{\bm x}f
+\frac{\bm F}{m}\cdot\nabla_{\bm v}f
=C[f].
\]
左侧是相空间输运，右侧是碰撞或松弛效应。

若忽略离散碰撞，带电粒子满足 Vlasov 方程：
\[
\frac{\partial f_s}{\partial t}
+\bm v\cdot\nabla_{\bm x}f_s
+\frac{q_s}{m_s}
(\bm E+\bm v\times\bm B)\cdot\nabla_{\bm v}f_s=0.
\]
电荷密度和电流密度由
\[
\rho_c=\sum_s q_s\int f_s\dd\bm v,
\qquad
\bm J=\sum_s q_s\int \bm v f_s\dd\bm v
\]
给出，再反馈到 Maxwell 方程中。由此形成 Vlasov--Maxwell 自洽系统。

对于稀薄中性气体，把碰撞近似为瞬时二体碰撞，可得玻尔兹曼方程
\[
\frac{\partial f}{\partial t}
+\bm v\cdot\nabla_{\bm x}f
+\frac{\bm F}{m}\cdot\nabla_{\bm v}f
=Q(f,f),
\]
其中
\[
Q(f,f)
=
\int B\left(f'f_*'-ff_*\right)\dd\Omega\dd\bm v_*.
\]
其关键封闭是假设碰撞前双粒子速度近似无关联：
\[
f_2(\bm v,\bm v_*)\approx f(\bm v)f(\bm v_*).
\]

完整碰撞积分计算昂贵，BGK 模型用单一松弛过程替代：
\[
C[f]=\frac{f^{\mathrm{eq}}-f}{\tau},
\]
其中
\[
f^{\mathrm{eq}}
=
n\left(\frac{m}{2\pi k_BT}\right)^{3/2}
\exp\left[
-\frac{m|\bm v-\bm u|^2}{2k_BT}
\right].
\]

对于等离子体中的长程 Coulomb 相互作用，粒子常经历大量小角度偏转。此时可得到 Landau 或 Fokker--Planck 型算子，把碰撞理解为速度空间中的漂移和扩散。

因此，碰撞项的不同假设对应一组模型层级：
\[
\text{无碰撞 Vlasov}
\longrightarrow
\text{弱碰撞 Vlasov--Fokker--Planck}
\longrightarrow
\text{强碰撞流体模型}.
\]

\chapter{从速度矩到连续介质方程}

分布函数包含完整速度分布，但宏观工程只保留少数矩。

数密度、质量密度和宏观速度分别为
\[
n=\int f\dd\bm v,
\qquad
\rho=mn,
\qquad
\rho\bm u=m\int \bm v f\dd\bm v.
\]

定义相对速度
\[
\bm c=\bm v-\bm u.
\]
压力张量与热流为
\[
\bm P=m\int \bm c\bm c f\dd\bm v,
\qquad
\bm q=\frac{m}{2}\int |\bm c|^2\bm c f\dd\bm v.
\]

对动理学方程取零阶矩，利用碰撞守恒粒子数：
\[
\int Q(f,f)\dd\bm v=0,
\]
得到连续性方程
\[
\boxed{
\frac{\partial\rho}{\partial t}
+\nabla\cdot(\rho\bm u)=0
}.
\]

其控制体形式为
\[
\frac{\dd}{\dd t}\int_\Omega \rho\dd V
=
-\int_{\partial\Omega}\rho\bm u\cdot\bm n\dd S.
\]
有限体积法正是从这一形式出发。

取一阶矩得到动量方程
\[
\frac{\partial(\rho\bm u)}{\partial t}
+
\nabla\cdot(\rho\bm u\bm u+\bm P)
=
\rho\bm b.
\]
写成
\[
\bm P=p\bm I-\bm\tau,
\]
则
\[
\frac{\partial(\rho\bm u)}{\partial t}
+
\nabla\cdot(\rho\bm u\bm u)
=
-\nabla p+\nabla\cdot\bm\tau+\rho\bm b.
\]

要封闭方程，必须给出应力与宏观变量之间的本构关系。Newton 流体满足
\[
\bm\tau
=
\mu\left[\nabla\bm u+(\nabla\bm u)^T\right]
+\lambda(\nabla\cdot\bm u)\bm I.
\]

不可压缩、常密度、常黏度情况下，
\[
\nabla\cdot\bm u=0,
\]
\[
\rho\left(
\frac{\partial\bm u}{\partial t}
+\bm u\cdot\nabla\bm u
\right)
=
-\nabla p+\mu\nabla^2\bm u+\rho\bm b.
\]

其中惯性项负责搬运动量，压力项负责重新分配动量，黏性项负责平滑速度梯度。若忽略黏性和热传导，则得到 Euler 方程。

取能量矩可得总能量方程
\[
\frac{\partial(\rho E)}{\partial t}
+
\nabla\cdot[(\rho E+p)\bm u]
=
\nabla\cdot(\bm\tau\cdot\bm u)
-\nabla\cdot\bm q
+\rho\bm b\cdot\bm u,
\]
其中
\[
E=e+\frac12|\bm u|^2,
\qquad
\bm q=-k\nabla T.
\]

从动理学到连续介质的本质，是只保留有限个矩，并用本构关系封闭更高阶矩。

\chapter{Maxwell 方程、线性化与微扰算子}

Maxwell 方程组为
\[
\nabla\cdot\bm E=\frac{\rho_c}{\varepsilon_0},
\qquad
\nabla\cdot\bm B=0,
\]
\[
\frac{\partial\bm B}{\partial t}+\nabla\times\bm E=0,
\]
\[
\frac{\partial\bm E}{\partial t}
-c^2\nabla\times\bm B
=
-\frac{\bm J}{\varepsilon_0}.
\]

对 Amp\`ere--Maxwell 方程取散度，并使用 Gauss 定律，得到
\[
\frac{\partial\rho_c}{\partial t}
+\nabla\cdot\bm J=0.
\]
因此 Maxwell 方程的连续结构自动包含电荷守恒。数值离散若破坏这一结构，可能产生非物理电荷。

带电连续介质中的 Lorentz 力为
\[
\bm f_L=\rho_c\bm E+\bm J\times\bm B.
\]
结合流体方程和 Ohm 定律，可得到双流体模型或单流体 MHD。

真空 Maxwell 方程是线性的，但 Maxwell 方程与介质、载流子、流体和非线性本构关系耦合后，整体系统通常是非线性的。设耦合模型抽象写成
\[
\Nop(\bm U)=0.
\]
在基态 \(\bm U_0\) 附近写
\[
\bm U=\bm U_0+\epsilon\bm u,
\qquad 0<\epsilon\ll1.
\]
若算子足够光滑，则
\[
\Nop(\bm U_0+\epsilon\bm u)
=
\Nop(\bm U_0)
+\epsilon D\Nop[\bm U_0](\bm u)
+\frac{\epsilon^2}{2}
D^2\Nop[\bm U_0](\bm u,\bm u)
+\cdots.
\]
若基态满足
\[
\Nop(\bm U_0)=0,
\]
保留一阶项得
\[
\Lop\bm u=0,
\qquad
\Lop:=D\Nop[\bm U_0].
\]

这里 \(D\Nop[\bm U_0]\) 是 Fr\'echet 导数，也可称线性化算子、Jacobian 算子或切线算子。它是非线性算子在基态处的局部线性近似。

例如非线性介质满足
\[
\bm D=\bm D(\bm E),
\qquad
\bm B=\bm B(\bm H).
\]
令
\[
\bm E=\bm E_0+\epsilon\bm e,
\qquad
\bm H=\bm H_0+\epsilon\bm h.
\]
本构关系的一阶展开为
\[
\bm D(\bm E_0+\epsilon\bm e)
=
\bm D(\bm E_0)
+
\epsilon
\left.
\frac{\partial\bm D}{\partial\bm E}
\right|_{\bm E_0}\bm e
+\Order(\epsilon^2),
\]
\[
\bm B(\bm H_0+\epsilon\bm h)
=
\bm B(\bm H_0)
+
\epsilon
\left.
\frac{\partial\bm B}{\partial\bm H}
\right|_{\bm H_0}\bm h
+\Order(\epsilon^2).
\]

定义切线材料张量
\[
\bm\varepsilon_{\mathrm{tan}}
=
\left.\frac{\partial\bm D}{\partial\bm E}\right|_{\bm E_0},
\qquad
\bm\mu_{\mathrm{tan}}
=
\left.\frac{\partial\bm B}{\partial\bm H}\right|_{\bm H_0}.
\]
若基态满足原方程，则一阶扰动满足线性 Maxwell 型系统
\[
\nabla\times\bm e
=
-\bm\mu_{\mathrm{tan}}\frac{\partial\bm h}{\partial t},
\]
\[
\nabla\times\bm h
=
\bm\varepsilon_{\mathrm{tan}}\frac{\partial\bm e}{\partial t}.
\]

在 Vlasov--Maxwell 系统中，设
\[
f_s=f_{s0}+\epsilon f_{s1},
\qquad
\bm E=\epsilon\bm E_1,
\qquad
\bm B=\bm B_0+\epsilon\bm B_1.
\]
忽略 \(\Order(\epsilon^2)\) 项，可得
\[
\frac{\partial f_{s1}}{\partial t}
+\bm v\cdot\nabla_{\bm x}f_{s1}
+\frac{q_s}{m_s}
(\bm v\times\bm B_0)\cdot\nabla_{\bm v}f_{s1}
=
-\frac{q_s}{m_s}
(\bm E_1+\bm v\times\bm B_1)
\cdot\nabla_{\bm v}f_{s0}.
\]
扰动电流为
\[
\bm J_1=\sum_s q_s\int \bm v f_{s1}\dd\bm v.
\]

若进一步假设平面波
\[
\bm u_1(\bm x,t)
=
\widehat{\bm u}
\exp[\ii(\bm k\cdot\bm x-\omega t)],
\]
则
\[
\partial_t\mapsto-\ii\omega,
\qquad
\nabla\mapsto\ii\bm k.
\]
线性偏微分方程组变为
\[
\bm M(\omega,\bm k)\widehat{\bm u}=0.
\]
存在非零解的条件
\[
\det\bm M(\omega,\bm k)=0
\]
就是色散关系。

微扰法、线性稳定性分析和 Newton 法共享同一局部线性思想，但用途不同：微扰法按小参数逐阶构造解，稳定性分析考察扰动增长率，Newton 法则用 Jacobian 求非线性残差的修正量。

\chapter{尺度分析与渐近展开}

常见无量纲数包括
\[
\mathrm{Kn}=\frac{\lambda}{L},
\qquad
\mathrm{Re}=\frac{\rho UL}{\mu},
\qquad
\mathrm{Ma}=\frac{U}{c}.
\]
它们不是标签，而是不同物理机制的相对量级。模型选择本质上是决定哪些项在给定极限下保留。

Taylor 展开与渐近展开需要严格区分。Taylor 展开通常由局部导数给出：
\[
f(x+h)
=
f(x)+hf'(x)+\frac{h^2}{2}f''(x)+\cdots.
\]
它特别适合分析离散算子与连续算子的局部差异。

渐近展开更一般。设 \(\epsilon\to0\)，若尺度函数满足
\[
\phi_{n+1}(\epsilon)=o(\phi_n(\epsilon)),
\]
且对每个固定 \(N\) 都有
\[
f(\epsilon)
-
\sum_{n=0}^{N}a_n\phi_n(\epsilon)
=
o(\phi_N(\epsilon)),
\]
则写作
\[
f(\epsilon)\sim\sum_{n=0}^{\infty}a_n\phi_n(\epsilon).
\]

渐近级数不要求无限求和收敛，只要求每次截断后，余项比最后保留的尺度更小。尺度也不必是整数幂，可以出现
\[
1,\quad
\epsilon^{1/2},\quad
\epsilon\log\epsilon,\quad
\ee^{-1/\epsilon}.
\]

因此，Taylor 级数是由局部导数产生、以普通幂为基的一类特殊展开；渐近展开强调的是不同项的尺度排序。

若系数增长很快，渐近级数可能先改善后恶化。固定 \(\epsilon\) 后，项
\[
|a_n\epsilon^n|
\]
可能先减小再增大。此时应在最小项附近截断，而不是追求无限求和。这称为最优截断。

若
\[
u=u_0+\epsilon u_1+\epsilon^2u_2+\cdots
\]
代入后，各阶问题保持原方程的微分阶数和边界条件结构，称为正则摄动。

若小参数乘在最高阶导数前，例如
\[
\epsilon u''+u'=0,
\]
则 \(\epsilon\to0\) 时方程阶数下降，部分边界条件无法由外部解满足。这是奇异摄动，通常需要边界层、匹配渐近展开、多尺度法或 WKB 方法。

\section{等离子体示例：小 Debye 长度与准中性极限}

考虑一维无量纲化的 Poisson 方程
\[
-\epsilon^2\frac{\partial^2\phi}{\partial x^2}
=
n_i-n_e,
\qquad
\epsilon=\frac{\lambda_D}{L}\ll1.
\]
这里 \(\lambda_D\) 是 Debye 长度，\(L\) 是宏观尺度。

若直接令 \(\epsilon=0\)，得到
\[
n_i=n_e,
\]
即准中性条件。这说明在主体区域，净电荷密度是高阶小量。

但如果边界给定的电势与主体准中性解不相容，最高阶导数前的小参数使问题成为奇异摄动。主体区域的外部展开可写成
\[
\phi^{\mathrm{out}}
=
\phi_0+\epsilon\phi_1+\epsilon^2\phi_2+\cdots,
\]
而边界附近必须引入拉伸坐标
\[
\xi=\frac{x}{\epsilon}.
\]
在该尺度下
\[
\frac{\partial}{\partial x}
=
\frac{1}{\epsilon}\frac{\partial}{\partial\xi},
\]
Poisson 方程恢复为
\[
-\frac{\partial^2\phi^{\mathrm{in}}}{\partial\xi^2}
=
n_i-n_e.
\]

这说明一个厚度为 \(\Order(\lambda_D)\) 的 Debye 鞘层可以在宏观尺度上极薄，却不能简单删除，因为它承担了满足壁面电势和电荷分离的任务。

该例展示了渐近展开的几个典型特征：
\begin{enumerate}
  \item 小参数来自明确的尺度比 \(\lambda_D/L\)；
  \item 主体区域和边界层使用不同坐标；
  \item 令小参数为零会降低方程阶数；
  \item 内外解必须在重叠区域匹配；
  \item 极限模型的有效性依赖于空间区域和边界条件。
\end{enumerate}

在等离子体文献中，类似结构还出现在低频极限、强磁场极限、漂移动理学、回旋动理学和 MHD 极限中。阅读时应先问：小参数是什么，哪些变量被假设为同阶，极限是否降低方程阶数，是否存在快时间、短空间或边界层尺度。

常见记号应按以下方式理解：
\[
f=g+\Order(\epsilon^p)
\]
表示余项被常数倍 \(\epsilon^p\) 控制；
\[
f=g+o(\epsilon^p)
\]
表示余项与 \(\epsilon^p\) 的比值趋于零；
\[
f\sim g
\]
既可能表示 \(f/g\to1\)，也可能表示逐阶渐近展开，必须结合上下文判断。

% =========================================================
\part{数值离散}

\chapter{一致性、稳定性与修正方程}

考虑线性平流方程
\[
u_t+au_x=0,
\qquad
u(x,t)=u_0(x-at).
\]

前向时间差分满足
\[
\frac{u_i^{n+1}-u_i^n}{\Delta t}
=
u_t+\Order(\Delta t),
\]
中心空间差分满足
\[
\frac{u_{i+1}-u_{i-1}}{2\Delta x}
=
u_x+\Order(\Delta x^2).
\]
离散格式的阶数来自离散算子与连续算子的局部展开匹配。

对线性常系数格式，可令误差模态为
\[
e_i^n=\widehat e^n\ee^{\ii kx_i},
\]
得到
\[
\widehat e^{n+1}=G(k)\widehat e^n.
\]
若存在波数使
\[
|G(k)|>1,
\]
则误差会增长。von Neumann 稳定性条件通常写成
\[
|G(k)|\le1.
\]

显式格式还受到 CFL 条件约束。定义
\[
\nu=\frac{a\Delta t}{\Delta x}.
\]
其含义是数值依赖域必须覆盖物理依赖域。

当 \(a>0\) 时，一阶上风格式为
\[
u_i^{n+1}
=
u_i^n-\nu(u_i^n-u_{i-1}^n)
=
(1-\nu)u_i^n+\nu u_{i-1}^n.
\]
在
\[
0\le\nu\le1
\]
时，新值是旧值的凸组合，因此不会产生新的局部极值。

其修正方程为
\[
u_t+au_x
=
\frac{a\Delta x}{2}(1-\nu)u_{xx}
+\cdots.
\]
右侧等效扩散项解释了上风格式的数值耗散：它抑制高频误差，也抹平真实小尺度。

\chapter{Godunov 定理与高分辨率格式}

Godunov 定理是守恒律数值计算中的基本限制。对线性标量平流问题，任何满足以下条件的线性格式：

\begin{enumerate}
  \item 线性；
  \item 单调，即不产生新的极值；
  \item 一致；
\end{enumerate}

其精度至多为一阶。

直观地说，高阶线性重构需要使用正负交替的权重来提高局部逼近阶数，而单调性要求更新具有凸组合性质。两者一般不能同时成立。

因此，不存在同时满足“线性、高阶、无振荡”的万能格式。Godunov 定理不是说高阶格式不可能，而是说必须放弃至少一个条件。现代高分辨率方法通常采取以下路线：

\begin{itemize}
  \item 保留守恒性；
  \item 允许算法是非线性的；
  \item 在光滑区域恢复高阶；
  \item 在间断附近退化为一阶单调格式。
\end{itemize}

以有限体积法为例，
\[
\frac{\dd \bar U_i}{\dd t}
=
-\frac{1}{\Delta x}
\left(
\widehat F_{i+1/2}
-
\widehat F_{i-1/2}
\right).
\]
关键在于构造界面左右状态
\[
U_{i+1/2}^{L},
\qquad
U_{i+1/2}^{R},
\]
并通过 Riemann 求解器或近似 Riemann 求解器计算数值通量。

MUSCL 方法采用分片线性重构，例如
\[
U_{i+1/2}^{L}
=
U_i+\frac12\sigma_i,
\]
其中斜率 \(\sigma_i\) 由限制器给出。典型限制器包括 minmod、van Leer 和 Superbee。限制器根据局部光滑性非线性地调整斜率，使格式在光滑区近似二阶，在激波或间断附近退化为一阶。

TVD 方法要求离散总变差不增加：
\[
TV(U^{n+1})\le TV(U^n),
\qquad
TV(U)=\sum_i|U_{i+1}-U_i|.
\]
它把“不过度产生振荡”转化为可分析的离散条件。

ENO/WENO 方法不直接限制单一斜率，而是在多个候选模板之间选择或加权。WENO 在光滑区域使用高阶组合，在包含间断的模板上自动降低权重。

Godunov 定理揭示的是一个结构性折中：
\[
\text{高阶精度}
\quad\leftrightarrow\quad
\text{单调性与无振荡}.
\]
工程格式设计的重点不是消除这个矛盾，而是让格式依据局部解结构自适应地处理这个矛盾。

\chapter{色散关系与波动误差}

设连续波解为
\[
u=\widehat u\ee^{\ii(kx-\omega t)}.
\]
对线性平流方程，
\[
\omega=ak.
\]
相速度和群速度均为 \(a\)：
\[
c_p=\frac{\omega}{k}=a,
\qquad
c_g=\frac{\dd\omega}{\dd k}=a.
\]

离散格式通常得到数值频率
\[
\omega_{\mathrm{num}}(k).
\]
若
\[
\operatorname{Re}\omega_{\mathrm{num}}\neq ak,
\]
则出现数值色散；若
\[
\operatorname{Im}\omega_{\mathrm{num}}<0,
\]
则出现数值耗散。

流体激波计算通常需要适量耗散，声学、电磁波和弹性波计算则更关注长距离传播中的相位误差。格式选择必须匹配目标物理。

\chapter{有限差分、有限体积、有限元与谱方法}

有限差分法直接近似导数，规则网格上高效，适合高阶格式和 Fourier 分析，但复杂几何处理困难。

有限体积法从控制体守恒式出发：
\[
\frac{\dd}{\dd t}
\int_{\Omega_i}U\dd V
+
\int_{\partial\Omega_i}
\bm F\cdot\bm n\dd S
=
\int_{\Omega_i}S\dd V.
\]
相邻单元共享界面通量，因此局部和全局守恒自然成立。

有限元法从弱形式出发。以 Poisson 方程
\[
-\nabla^2u=f
\]
为例，乘测试函数并分部积分：
\[
\int_\Omega\nabla v\cdot\nabla u\dd V
=
\int_\Omega vf\dd V+\text{边界项}.
\]
展开
\[
u_h=\sum_jU_j\phi_j
\]
后得到
\[
\bm K\bm U=\bm F.
\]

谱方法使用 Fourier 或 Chebyshev 等全局基函数，对光滑解可获得很高精度，但复杂几何、局部加密和间断处理更困难。

% =========================================================
\part{矩阵求解}

\chapter{从离散算子到代数系统}

偏微分方程离散后通常得到
\[
\bm A\bm x=\bm b
\]
或非线性残差系统
\[
\bm R(\bm U)=0.
\]

局部微分算子通常产生稀疏矩阵，边界积分算子常产生稠密矩阵。矩阵结构直接反映物理耦合范围。

对于非线性系统，Newton 法在第 \(k\) 步求解
\[
\bm J(\bm U^k)\Delta\bm U
=
-\bm R(\bm U^k),
\qquad
\bm U^{k+1}=\bm U^k+\Delta\bm U.
\]
大型非线性计算通常是外层 Newton 或伪时间迭代，加内层稀疏线性求解。

\chapter{LU 与 LU-SGS}

稠密直接法可将
\[
\bm A=\bm L\bm U,
\]
然后依次求解
\[
\bm L\bm y=\bm b,
\qquad
\bm U\bm x=\bm y.
\]
稠密 LU 的分解复杂度约为
\[
O(N^3),
\]
存储量为
\[
O(N^2).
\]

稀疏矩阵做精确 LU 时会产生填充，因此大型三维 PDE 通常需要迭代法或不完全分解。

LU-SGS 将
\[
\bm A
\]
近似分解为
\[
\bm A
\approx
(\bm D+\bm L)\bm D^{-1}(\bm D+\bm U),
\]
通过前向和后向扫掠近似求解。它常用于可压缩 CFD 的隐式推进，也可作为预条件器。

\chapter{Krylov 方法、Arnoldi、GMRES 与 BiCGSTAB}

给定初始残差
\[
\bm r_0=\bm b-\bm A\bm x_0,
\]
Krylov 子空间为
\[
\Ksp_m(\bm A,\bm r_0)
=
\operatorname{span}
\{
\bm r_0,\bm A\bm r_0,\dots,\bm A^{m-1}\bm r_0
\}.
\]

Arnoldi 过程构造正交基
\[
\bm V_m=[\bm v_1,\dots,\bm v_m],
\]
满足
\[
\bm A\bm V_m
=
\bm V_{m+1}\overline{\bm H}_m.
\]

GMRES 在
\[
\bm x_0+\Ksp_m
\]
中最小化残差二范数：
\[
\bm x_m
=
\arg\min_{\bm x\in\bm x_0+\Ksp_m}
\|\bm b-\bm A\bm x\|_2.
\]
其优点是对一般非对称系统较稳健，缺点是基向量存储和正交化成本随迭代步增长。重启 GMRES 通过限制子空间维数控制成本，但可能损失收敛信息。

BiCGSTAB 使用短递推，内存固定，单步成本较低，适用于非对称稀疏系统。它的收敛曲线可能振荡，并对预条件质量较敏感。

\chapter{预条件与 AMG}

条件数
\[
\kappa(\bm A)=\|\bm A\|\,\|\bm A^{-1}\|
\]
衡量扰动放大和迭代难度。

预条件器 \(\bm M\) 的目标是使
\[
\bm M^{-1}\bm A
\]
具有更有利的谱性质，同时保证应用 \(\bm M^{-1}\) 的成本足够低。

局部平滑器擅长消除高频误差，但低频误差衰减慢。多重网格把细网格低频误差转移到粗网格，在粗层上它表现为较高频误差，因此更容易消除。

AMG 不依赖显式几何层级，而根据矩阵的强连接关系构造粗空间。典型循环包括：
\begin{enumerate}
  \item 预平滑；
  \item 计算残差；
  \item 限制到粗层；
  \item 求解粗层修正；
  \item 延拓回细层；
  \item 后平滑。
\end{enumerate}

AMG 对 Poisson、压力修正和扩散型系统尤其有效，也常作为 GMRES 的预条件器。

\chapter{MoM 与稠密积分方程}

计算电磁学中的矩量法通常从积分方程
\[
\Lop J=g
\]
出发。展开
\[
J=\sum_{n=1}^N I_nf_n,
\]
并用测试函数 \(w_m\) 投影：
\[
\avg{w_m,\Lop J}
=
\avg{w_m,g}.
\]
得到
\[
\bm Z\bm I=\bm V.
\]

由于积分核具有长程作用，\(\bm Z\) 通常是稠密矩阵。大规模 MoM 常与快速多极子、FFT 加速、层次矩阵和 Krylov 求解器结合。

% =========================================================
\part{网格技巧}

\chapter{网格是几何与代数之间的接口}

网格决定几何表示误差、局部离散误差、稳定性限制、矩阵稀疏结构和并行通信量。工业计算中，网格处理往往比方程推导更消耗工程时间。

贴体网格使网格边界与物体边界一致。其边界条件施加自然，适合解析壁面边界层，但复杂几何和大运动会导致网格生成与质量控制困难。

\chapter{浸没边界法}

浸没边界法把复杂物体放入简单背景网格，不要求网格贴合物体。在动量方程中加入边界作用：
\[
\rho\left(
\frac{\partial\bm u}{\partial t}
+\bm u\cdot\nabla\bm u
\right)
=
-\nabla p+\mu\nabla^2\bm u+\bm f_{\mathrm{IB}}.
\]

连续强迫法通过光滑核在拉格朗日边界与欧拉网格之间传递力；离散强迫法直接修改边界附近离散方程。

其优点是网格规则、复杂运动处理方便；缺点是壁面附近精度、严格守恒和薄结构处理较困难。

\chapter{四叉树、八叉树与自适应加密}

二维四叉树将一个单元递归分成四个子单元，三维八叉树分成八个子单元。其核心是只在激波、剪切层、涡、界面或高误差区域投入高分辨率。

常见工程约束包括：
\begin{itemize}
  \item 相邻单元保持 2:1 层级平衡；
  \item 处理粗细非匹配界面；
  \item 加密和粗化时守恒地转移变量；
  \item 自适应后重新做并行负载均衡。
\end{itemize}

\chapter{动网格与几何守恒律}

在任意拉格朗日--欧拉描述中，网格具有速度 \(\bm w\)。守恒方程写成
\[
\frac{\dd}{\dd t}\int_{\Omega(t)}U\dd V
+
\int_{\partial\Omega(t)}
[\bm F(U)-U\bm w]\cdot\bm n\dd S
=0.
\]

动网格必须满足几何守恒律：若物理解为常数，仅仅移动网格不应制造虚假变化。

网格运动可由弹簧类比、Laplace 平滑、线弹性、径向基函数或局部重构实现。运动过大时，需要重新划分网格。

\chapter{重叠网格}

重叠网格将多个独立网格块叠加，例如背景网格、贴体部件网格和随动局部网格。网格之间通过供体单元向接收单元插值。

其优势是部件网格可独立生成，适合多个刚体的大范围相对运动。困难在于孔切、供体搜索、守恒插值和动态连通更新。

\chapter{工程选择原则}

若几何固定且壁面精度要求高，优先使用贴体网格。

若几何复杂、运动明显，而壁面剪切应力不是首要目标，可使用浸没边界法。

若局部尺度变化极大，可使用四叉树或八叉树自适应网格。

若边界运动幅度较小，可使用动网格。

若多个刚体发生大幅相对运动，可使用重叠网格。

实际工业软件常组合使用：
\[
\text{贴体边界层}
+
\text{自适应远场},
\]
\[
\text{八叉树背景网格}
+
\text{浸没边界},
\]
\[
\text{重叠网格}
+
\text{局部动网格}.
\]

\chapter{完整计算工作流}

一次可靠计算通常按以下链条进行：

\begin{enumerate}
  \item 用 \(\mathrm{Kn}\)、\(\mathrm{Re}\)、\(\mathrm{Ma}\) 等尺度比选择模型；
  \item 确定守恒变量、材料模型和边界条件；
  \item 选择与几何和局部尺度匹配的网格策略；
  \item 选择满足守恒、稳定性和波动精度要求的离散格式；
  \item 根据显式或隐式推进形成代数系统；
  \item 根据矩阵对称性、稀疏性和谱性质选择求解器与预条件器；
  \item 分别监控模型误差、网格误差、离散误差、线性残差和非线性残差。
\end{enumerate}

残差下降只说明离散代数系统被求解，并不自动说明物理模型正确、网格充分或离散误差可接受。

\chapter{结语}

本书的统一主线是
\[
\boxed{
\text{刘维尔}
\rightarrow
\text{Vlasov / Boltzmann}
\rightarrow
\text{连续介质方程}
}
\]
\[
\boxed{
\text{渐近极限与线性化}
\rightarrow
\text{数值离散}
\rightarrow
\text{稳定性、耗散、色散}
}
\]
\[
\boxed{
\text{稀疏或稠密矩阵}
\rightarrow
\text{Krylov、LU、AMG、MoM}
\rightarrow
\text{工程网格}
}
\]

计算力学的本质，是先决定保留哪些物理信息，再决定如何离散这些信息，最后利用算子和矩阵结构把离散模型高效、稳定地计算出来。

\end{document}
